\documentclass[11pt]{article}
\usepackage{fontspec}
\directlua{luaotfload.add_fallback("hebfb",{"DejaVu Sans:mode=harf"})}
\usepackage{polyglossia}
\setdefaultlanguage{hebrew}
\setotherlanguage{english}
\setmainfont{Noto Sans Hebrew}[Script=Hebrew,RawFeature={fallback=hebfb}]
\newfontfamily\codefont{DejaVu Sans Mono}[RawFeature={fallback=hebfb}]
\usepackage[a4paper,margin=18mm]{geometry}
\usepackage[table]{xcolor}
\usepackage{mdframed}
\usepackage{tabularx}
\usepackage{xltabular}
\usepackage{array}
\usepackage{enumitem}
\usepackage{fancyhdr}
\setlist{leftmargin=1.6em,itemsep=1pt,topsep=2pt}
% colors
\definecolor{h1}{HTML}{1E3A8A}\definecolor{h2}{HTML}{1E40AF}\definecolor{h3}{HTML}{0F172A}
\definecolor{subgray}{HTML}{64748B}\definecolor{hdrbg}{HTML}{E0E7FF}
\definecolor{cfg}{HTML}{E2E8F0}\definecolor{cbg}{HTML}{0F172A}\definecolor{cbold}{HTML}{7DD3FC}\definecolor{ccom}{HTML}{94A3B8}
\definecolor{metabg}{HTML}{F1F5F9}
\newcommand{\enLR}[1]{\textenglish{#1}}
\newcommand{\code}[1]{\textenglish{\codefont #1}}
\newcommand{\chapterheading}[1]{\vspace{2pt}{\LARGE\bfseries\color{h1}#1\par}\vskip2pt{\color{h1}\hrule height 2pt}\vskip6pt}
\newcommand{\sectionheading}[1]{\vskip10pt\penalty-200{\large\bfseries\color{h2}#1\par}\vskip3pt}
\newcommand{\subheading}[1]{\vskip6pt{\bfseries\color{h3}#1\par}\vskip2pt}
\newcommand{\boxtitle}[1]{{\bfseries\color{boxti}#1}}
% box environments (right border, tinted bg)
\newmdenv[topline=false,bottomline=false,leftline=false,rightline=true,linewidth=2pt,innermargin=0pt,skipabove=6pt,skipbelow=6pt,innertopmargin=6pt,innerbottommargin=6pt,innerleftmargin=8pt,innerrightmargin=8pt]{genericbox}
\definecolor{faqbg}{HTML}{ECFDF5}\definecolor{faqfr}{HTML}{10B981}\definecolor{faqti}{HTML}{065F46}
\definecolor{misbg}{HTML}{FEF2F2}\definecolor{misfr}{HTML}{EF4444}\definecolor{misti}{HTML}{991B1B}
\definecolor{tipbg}{HTML}{FFFBEB}\definecolor{tipfr}{HTML}{F59E0B}\definecolor{tipti}{HTML}{92400E}
\definecolor{exabg}{HTML}{EFF6FF}\definecolor{exafr}{HTML}{3B82F6}\definecolor{exati}{HTML}{1E3A8A}
\definecolor{defbg}{HTML}{F5F3FF}\definecolor{deffr}{HTML}{8B5CF6}\definecolor{defti}{HTML}{5B21B6}
\definecolor{stobg}{HTML}{FDF4FF}\definecolor{stofr}{HTML}{D946EF}\definecolor{stoti}{HTML}{86198F}
\definecolor{exbg}{HTML}{F0FDFA}\definecolor{exfr}{HTML}{14B8A6}\definecolor{exti}{HTML}{115E59}
\newenvironment{faqbox}{\colorlet{boxti}{faqti}\begin{mdframed}[backgroundcolor=faqbg,linecolor=faqfr,rightline=true,leftline=false,topline=false,bottomline=false,linewidth=2pt,innertopmargin=6pt,innerbottommargin=6pt,innerleftmargin=8pt,innerrightmargin=8pt,skipabove=6pt,skipbelow=6pt]}{\end{mdframed}}
\newenvironment{mistakebox}{\colorlet{boxti}{misti}\begin{mdframed}[backgroundcolor=misbg,linecolor=misfr,rightline=true,leftline=false,topline=false,bottomline=false,linewidth=2pt,innertopmargin=6pt,innerbottommargin=6pt,innerleftmargin=8pt,innerrightmargin=8pt,skipabove=6pt,skipbelow=6pt]}{\end{mdframed}}
\newenvironment{tipbox}{\colorlet{boxti}{tipti}\begin{mdframed}[backgroundcolor=tipbg,linecolor=tipfr,rightline=true,leftline=false,topline=false,bottomline=false,linewidth=2pt,innertopmargin=6pt,innerbottommargin=6pt,innerleftmargin=8pt,innerrightmargin=8pt,skipabove=6pt,skipbelow=6pt]}{\end{mdframed}}
\newenvironment{exambox}{\colorlet{boxti}{exati}\begin{mdframed}[backgroundcolor=exabg,linecolor=exafr,rightline=true,leftline=false,topline=false,bottomline=false,linewidth=2pt,innertopmargin=6pt,innerbottommargin=6pt,innerleftmargin=8pt,innerrightmargin=8pt,skipabove=6pt,skipbelow=6pt]}{\end{mdframed}}
\newenvironment{defbox}{\colorlet{boxti}{defti}\begin{mdframed}[backgroundcolor=defbg,linecolor=deffr,rightline=true,leftline=false,topline=false,bottomline=false,linewidth=2pt,innertopmargin=6pt,innerbottommargin=6pt,innerleftmargin=8pt,innerrightmargin=8pt,skipabove=6pt,skipbelow=6pt]}{\end{mdframed}}
\newenvironment{storybox}{\colorlet{boxti}{stoti}\begin{mdframed}[backgroundcolor=stobg,linecolor=stofr,rightline=true,leftline=false,topline=false,bottomline=false,linewidth=2pt,innertopmargin=6pt,innerbottommargin=6pt,innerleftmargin=8pt,innerrightmargin=8pt,skipabove=6pt,skipbelow=6pt]}{\end{mdframed}}
\newenvironment{exbox}{\colorlet{boxti}{exti}\begin{mdframed}[backgroundcolor=exbg,linecolor=exfr,rightline=true,leftline=false,topline=false,bottomline=false,linewidth=2pt,innertopmargin=6pt,innerbottommargin=6pt,innerleftmargin=8pt,innerrightmargin=8pt,skipabove=6pt,skipbelow=6pt]}{\end{mdframed}}
\newenvironment{metabox}{\begin{mdframed}[backgroundcolor=metabg,linewidth=0pt,innertopmargin=5pt,innerbottommargin=5pt,innerleftmargin=8pt,innerrightmargin=8pt,skipabove=4pt,skipbelow=8pt]}{\end{mdframed}}
\newmdenv[backgroundcolor=cbg,linewidth=0pt,innertopmargin=6pt,innerbottommargin=6pt,innerleftmargin=8pt,innerrightmargin=8pt,skipabove=6pt,skipbelow=8pt]{codebox}
\setlength{\parindent}{0pt}\setlength{\parskip}{4pt}
\renewcommand{\thepage}{\enLR{\arabic{page}}}
\pagestyle{fancy}\fancyhf{}\fancyfoot[C]{\thepage}\renewcommand{\headrulewidth}{0pt}
\begin{document}
\sloppy
\chapterheading{פרק \enLR{1} – מבוא לאיומי רשת חדשים}{\small\color{subgray}\enLR{10} שעות עיוני · שבועות \enLR{1}–\enLR{3}\par}\vskip4pt

\begin{metabox}\textbf{מטרות:} אבולוציית אבטחת המידע · כלים לאבטחה · מערכות והתקפת מערכות\\ \textbf{בבגרות:} \enLR{phishing}, סוגי תוקפים, \enLR{brute force, DoS, buffer overflow}\end{metabox}

\begin{defbox}\boxtitle{לפני שמתחילים – מושגי יסוד ברשת (למי שאין רקע)}\par  כדי להבין איומים ואבטחה, צריך קודם להבין מה זו בכלל רשת. הנה כל מה שצריך, בשפה פשוטה:  \textbf{מחשב מארח \enLR{(Host)}} – כל מכשיר ברשת: מחשב, טלפון, שרת, מדפסת. \textbf{כתובת \enLR{IP}} – ה"כתובת" של המכשיר ברשת, כמו כתובת בית (למשל \enLR{192.168.1.10).} לפיה יודעים לאן לשלוח מידע. \textbf{כתובת \enLR{MAC}} – מספר סידורי ייחודי "צרוב" בכרטיס הרשת של כל מכשיר (למשל \enLR{0C:0E:15:22:05:97).} כמו מספר שלדה של רכב. \textbf{חבילה \enLR{(Packet)}} – המידע ברשת לא נשלח כמקשה אחת אלא נחתך ל"מעטפות" קטנות. כל חבילה נושאת כתובת מקור וכתובת יעד, בדיוק כמו מעטפה בדואר. \textbf{שרת \enLR{(Server)} ולקוח \enLR{(Client)}} – הלקוח מבקש שירות (הדפדפן שלכם), השרת נותן אותו (אתר האינטרנט). הלקוח שואל, השרת עונה. \textbf{נתב \enLR{(Router)}} – מכשיר שמעביר חבילות \emph{בין רשתות שונות} (למשל בין הרשת בבית לאינטרנט). "צומת דרכים". \textbf{מתג \enLR{(Switch)}} – מכשיר שמחבר מכשירים \emph{בתוך אותה רשת מקומית} (כל המחשבים במשרד). \textbf{פרוטוקול} – "שפה"/כללים מוסכמים שמכשירים מדברים בהם \enLR{(HTTP} לאתרים, \enLR{DNS} לתרגום שמות). כמו כללי נימוס בשיחת טלפון. \textbf{פורט \enLR{(Port)}} – "דלת" ממוספרת במכשיר לכל שירות \enLR{(80} לאתר, \enLR{443} לאתר מאובטח, \enLR{22} ל-\enLR{SSH).} לכתובת \enLR{IP} יש \enLR{65,535} דלתות כאלה.  כל אלה מוסברים בהרחבה ביחידת "תזכורת רשתות תקשורת". \textbf{הרעיון המרכזי של הקורס:} בכל נקודה כזו – מכשיר, כתובת, חבילה, פורט – תוקף יכול לנסות לתקוף, ואנחנו נלמד להגן.\end{defbox}

\sectionheading{\enLR{1.1} מהי אבטחת מידע ולמה הרשת היא שדה הקרב}

אבטחת מידע \enLR{(Information Security)} היא ההגנה על מידע ועל המערכות שמעבדות אותו מפני גישה, שימוש, חשיפה, שיבוש, שינוי או השמדה בלתי מורשים. \textbf{אבטחת רשת} \enLR{(Network Security)} היא תת-תחום: ההגנה על התשתית שמעבירה את המידע – נתבים, מתגים, חומות אש, קווי תקשורת – ועל התעבורה שעוברת בהם.\par

ההבחנה המרכזית שכדאי להציג לתלמידים כבר ביום הראשון: \textbf{נכס \enLR{(asset)} · איום \enLR{(threat)} · פגיעות \enLR{(vulnerability)} · סיכון \enLR{(risk)} · בקרה \enLR{(control)}}.\par

\begin{xltabular}{\linewidth}{|>{\setRTL\raggedleft\arraybackslash}X|>{\setRTL\raggedleft\arraybackslash}X|}
\hline
מה שמגנים עליו: מידע לקוחות, שרתים, מוניטין, זמינות שירות. & \textbf{נכס \enLR{(Asset)}} \\
\hline
חולשה במערכת: באג בתוכנה, סיסמה חלשה, פורט פתוח, עובד שלא הוכשר. & \textbf{פגיעות \enLR{(Vulnerability)}} \\
\hline
גורם שעלול לנצל פגיעות: האקר, נוזקה, עובד ממורמר, גם שריפה או הצפה. & \textbf{איום \enLR{(Threat)}} \\
\hline
הדרך/הקוד הקונקרטי שבו האיום מנצל את הפגיעות. & \textbf{ניצול \enLR{(Exploit)}} \\
\hline
הסתברות שהאיום יממש את הפגיעות × הנזק שייגרם. סיכון = איום × פגיעות × ערך הנכס. & \textbf{סיכון \enLR{(Risk)}} \\
\hline
אמצעי שמקטין סיכון: חומת אש, הצפנה, נוהל, מנעול. & \textbf{בקרה \enLR{(Control / Countermeasure)}} \\
\hline
\end{xltabular}

\subheading{שלושת עקרונות היסוד – \enLR{CIA Triad}}

\begin{xltabular}{\linewidth}{|>{\setRTL\raggedleft\arraybackslash}X|>{\setRTL\raggedleft\arraybackslash}X|>{\setRTL\raggedleft\arraybackslash}X|>{\setRTL\raggedleft\arraybackslash}X|}
\hline
\rowcolor{hdrbg}\textbf{אמצעי הגנה אופייני} & \textbf{דוגמה לפגיעה} & \textbf{משמעות} & \textbf{עיקרון} \\
\hline
\endhead
הצפנה \enLR{(AES, SSH, VPN)}, בקרת גישה, \enLR{ACL} & האזנה \enLR{(sniffing)} ל-\enLR{Telnet}, דליפת \enLR{DB} & רק מי שמורשה רואה את המידע & \textbf{סודיות} \enLR{(Confidentiality)} \\
\hline
גיבוב \enLR{(SHA), HMAC}, חתימה דיגיטלית & שינוי סכום בהעברה בנקאית, \enLR{MITM} & המידע לא שונה בדרך או במנוחה ללא הרשאה & \textbf{שלמות} \enLR{(Integrity)} \\
\hline
יתירות, גיבויים, \enLR{IPS}, סינון תעבורה & \enLR{DDoS}, כופרה, כשל חומרה & המידע והשירות נגישים כשצריך & \textbf{זמינות} \enLR{(Availability)} \\
\hline
\end{xltabular}

שני מושגים נלווים שנשאלים לעיתים: \textbf{אימות} \enLR{(Authentication} – מי אתה) ו\textbf{אי-הכחשה} \enLR{(Non-repudiation} – השולח לא יכול להכחיש ששלח; מושג באמצעות חתימה דיגיטלית, פרק \enLR{7).}\par

\begin{faqbox}\boxtitle{שאלת תלמיד: "למה צריך שלושה? הרי אם מצפינים הכול – הכול בסדר"}\par  הצפנה נותנת סודיות בלבד. קובץ מוצפן עדיין יכול להימחק (פגיעה בזמינות) או להיות מוחלף בקובץ מוצפן אחר (שלמות). כופרה \enLR{(Ransomware)} היא הדוגמה המושלמת: היא \emph{מצפינה} את הקבצים שלכם – ובכך פוגעת בזמינות. הצפנה בידיים הלא נכונות היא נשק.\end{faqbox}

\begin{storybox}\boxtitle{סיפור מהחיים: בית החולים שביטל ניתוחים \enLR{(WannaCry}, מאי \enLR{2017)}}\par  בבוקר יום שישי אחד, מחשבים ב-\enLR{80} בתי חולים של ה-\enLR{NHS} באנגליה הציגו מסך אדום: "הקבצים שלך הוצפנו. שלם \enLR{300}\$ בביטקוין". המערכות לתורים, לתוצאות בדיקות ולמכשירי הדמיה הפסיקו לעבוד. \enLR{19,000} תורים וניתוחים בוטלו, אמבולנסים הופנו לבתי חולים אחרים. הנוזקה לא "גנבה" כלום – היא פשוט מנעה גישה. זו פגיעה ב\textbf{זמינות}. ואיך היא התפשטה? דרך פגיעות ב-\enLR{SMB} של \enLR{Windows} שמיקרוסופט \emph{כבר תיקנה חודשיים קודם} – אבל בתי החולים לא התקינו את העדכון. המסקנה לכיתה: עדכון תוכנה הוא אמצעי אבטחה, ולא מטרד.\end{storybox}

\sectionheading{\enLR{1.2} אבולוציית אבטחת הרשת – ציר זמן}

\begin{xltabular}{\linewidth}{|>{\setRTL\raggedleft\arraybackslash}X|>{\setRTL\raggedleft\arraybackslash}X|>{\setRTL\raggedleft\arraybackslash}X|}
\hline
\rowcolor{hdrbg}\textbf{מה השתנה בעקבותיו} & \textbf{אירוע} & \textbf{שנה} \\
\hline
\endhead
ההוכחה שקוד יכול להתפשט ברשת & \enLR{Creeper} – "וירוס" הניסוי הראשון ב-\enLR{ARPANET}, ואחריו \enLR{Reaper} שנכתב כדי למחוק אותו & \enLR{1971} \\
\hline
תחילת המחקר האקדמי & פרד כהן טובע את המונח "וירוס מחשב" & \enLR{1983} \\
\hline
הקמת ה-\enLR{CERT} הראשון \enLR{(Carnegie Mellon)}; החוק הפדרלי הראשון להרשעה בפריצה & \textbf{תולעת מוריס} – \textasciitilde{}\enLR{10\%} ממחשבי האינטרנט מושבתים תוך יום. ניצלה חולשות ב-\enLR{sendmail} וב-\enLR{finger} וניחוש סיסמאות & \enLR{1988} \\
\hline
מודל "הגנה היקפית" – פנים מהימן, חוץ עוין & חומות אש מסחריות \enLR{(Check Point FireWall-1, 1994), IDS} ראשונים \enLR{(Snort, 1998), SSL (1995)} & \enLR{1990s} \\
\hline
תולעים שמנצלות חולשות ידועות ⇒ תרבות ה-\enLR{patching, Windows Update} אוטומטי, \enLR{IPS} & \enLR{ILOVEYOU (2000), Code Red (2001), SQL Slammer (2003} – \enLR{75,000} שרתים ב-\enLR{10} דקות), \enLR{Blaster, Sasser} & \enLR{2000}–\enLR{2004} \\
\hline
לוחמת סייבר בין מדינות; \enLR{APT} (איום מתמשך מתקדם) & מתקפת \enLR{DDoS} על אסטוניה \enLR{(2007)}; \textbf{\enLR{Stuxnet}} \enLR{(2010)} – נוזקה שפגעה בצנטריפוגות גרעיניות & \enLR{2007}–\enLR{2010} \\
\hline
אבטחת שרשרת אספקה, \enLR{IoT}, כופרה כמודל עסקי & פריצת \enLR{Target (40} מיליון כרטיסי אשראי דרך ספק מיזוג אוויר), \enLR{Mirai botnet (2016} – מצלמות \enLR{IoT)}, \textbf{\enLR{WannaCry}} \enLR{(2017} – כופרה שניצלה \enLR{SMB, 200,000} מחשבים ב-\enLR{150} מדינות) & \enLR{2013}–\enLR{2017} \\
\hline
מודל \textbf{\enLR{Zero Trust}} – "לעולם אל תסמוך, תמיד תאמת"; אין יותר "פנים מהימן" & \enLR{SolarWinds (2020} – פריצה דרך עדכון תוכנה), \enLR{Log4Shell (2021)}, כופרה כשירות \enLR{(RaaS)}, דיפ-פייק והנדסה חברתית מבוססת \enLR{AI} & \enLR{2020}–היום \\
\hline
\end{xltabular}

\begin{tipbox}\boxtitle{טיפ להוראה}\par  אל תקריאו את הטבלה. בחרו \enLR{3} סיפורים ותספרו אותם: תולעת מוריס (סטודנט שרצה "למדוד את האינטרנט" והסתבך), \enLR{Target} (איך מזגן הפיל רשת ענק) ו-\enLR{WannaCry} (בית חולים באנגליה שביטל ניתוחים). תלמידים זוכרים סיפורים, לא שנים.\end{tipbox}

\begin{storybox}\boxtitle{סיפור מהחיים: הסטודנט שהפיל את האינטרנט בטעות (תולעת מוריס, \enLR{1988)}}\par  רוברט מוריס, סטודנט בן \enLR{23} בקורנל, כתב תוכנית שנועדה "למדוד כמה גדול האינטרנט" – היא הייתה אמורה להעתיק את עצמה למחשב פעם אחת. בגלל באג, היא שאלה כל מחשב "אתה כבר נגוע?" ואם הוא ענה "כן" – בכל זאת הדביקה אותו מחדש ב-\enLR{1} מ-\enLR{7} מהמקרים (מוריס חשש שמנהלים יזייפו "כן"). התוצאה: כל מחשב הריץ מאות עותקים עד שקרס. \enLR{6,000} מחשבים – כעשירית מהאינטרנט דאז – הושבתו. מוריס היה האדם הראשון שהורשע לפי חוק הונאת המחשבים האמריקאי. היום הוא פרופסור ב-\enLR{MIT.} השאלה לכיתה: האם זה היה וירוס או תולעת? (תולעת – התפשטה לבד דרך הרשת ללא מגע אדם.)\end{storybox}

\sectionheading{\enLR{1.3} מי תוקף? סוגי תוקפים \enLR{(Threat Actors)}}

נושא שנשאל ישירות בבגרות (שאלה \enLR{6}ג). חשוב להבחין לפי \textbf{מניע} ו\textbf{יכולת}:\par

\begin{xltabular}{\linewidth}{|>{\setRTL\raggedleft\arraybackslash}X|>{\setRTL\raggedleft\arraybackslash}X|>{\setRTL\raggedleft\arraybackslash}X|>{\setRTL\raggedleft\arraybackslash}X|}
\hline
\rowcolor{hdrbg}\textbf{רמת יכולת} & \textbf{מניע} & \textbf{מי הוא} & \textbf{סוג} \\
\hline
\endhead
נמוכה – אבל הכלים חזקים & סקרנות, יוקרה, שעמום & חובבן שמריץ כלים שאחרים כתבו, בלי להבין את התהליך & \textbf{\enLR{Script Kiddie}} \\
\hline
בינונית & מחאה, השחתת אתרים, הדלפות & תוקף ממניעים אידיאולוגיים/פוליטיים \enLR{(Anonymous)} & \textbf{\enLR{Hacktivist}} \\
\hline
גבוהה, משאבים רבים & כסף: כופרה, גניבת כרטיסי אשראי, הונאות & קבוצה מאורגנת שמטרתה רווח כספי & \textbf{פשע מאורגן} \enLR{(Organized crime / Cybercriminal)} \\
\hline
מסוכן במיוחד – כבר בפנים & נקמה, כסף, טעות & עובד/קבלן עם גישה מורשית שמשתמש בה לרעה (או ברשלנות) & \textbf{\enLR{Insider}} (איום פנימי) \\
\hline
הגבוהה ביותר, סבלנות של שנים & ריגול, חבלה בתשתיות & יחידות סייבר ממשלתיות & \textbf{מדינה} \enLR{(State-sponsored / APT)} \\
\hline
 &  & חלוקה לפי אתיקה: \enLR{white} – בודק חדירה באישור; \enLR{black} – פושע; \enLR{gray} – פורץ ללא אישור אבל מדווח & \textbf{\enLR{White hat / Gray hat / Black hat}} \\
\hline
\end{xltabular}

\begin{storybox}\boxtitle{סיפור מהחיים: איך מזגן הפיל רשת קמעונאית \enLR{(Target, 2013)}}\par  רשת החנויות \enLR{Target} איבדה פרטי \enLR{40} מיליון כרטיסי אשראי. התוקפים לא פרצו ל-\enLR{Target} ישירות – הם שלחו \enLR{phishing} לחברה קטנה בפנסילבניה שמתחזקת את מערכות המיזוג של \enLR{Target}, גנבו את פרטי ההתחברות שלה לפורטל הספקים, ומשם – כי הרשת לא הייתה מפולחת – הגיעו לקופות. שלושה לקחים שממפים ישירות לקורס: \enLR{(1)} הנדסה חברתית היא נקודת הכניסה הנפוצה ביותר; \enLR{(2) "}יחסי אמון" עם צד שלישי הם פגיעות \enLR{(Trust exploitation); (3)} פילוח רשת \enLR{(VLAN, ACL} – פרקים \enLR{4} ו-\enLR{6)} היה עוצר את התנועה הרוחבית.\end{storybox}

\sectionheading{\enLR{1.4} נוזקות: וירוסים, תולעים, סוסים טרויאנים ועוד}

\textbf{נוזקה} \enLR{(Malware = malicious software)} הוא שם כולל. ההבחנה בין שלושת הסוגים הקלאסיים היא לפי \emph{איך הם מתפשטים}:\par

\begin{xltabular}{\linewidth}{|>{\setRTL\raggedleft\arraybackslash}X|>{\setRTL\raggedleft\arraybackslash}X|>{\setRTL\raggedleft\arraybackslash}X|>{\setRTL\raggedleft\arraybackslash}X|}
\hline
\rowcolor{hdrbg}\textbf{סוס טרויאני \enLR{(Trojan)}} & \textbf{תולעת \enLR{(Worm)}} & \textbf{וירוס \enLR{(Virus)}} & \textbf{} \\
\hline
\endhead
תוכנה שנראית לגיטימית אך מכילה קוד זדוני & תוכנית עצמאית שמשכפלת את עצמה דרך הרשת & קוד זדוני שמצמיד את עצמו לקובץ/תוכנה לגיטימית & \textbf{הגדרה} \\
\hline
הוא עצמו "המארח" המזויף & לא – עומד בפני עצמו & כן – נדבק לקובץ & \textbf{צריך מארח?} \\
\hline
המשתמש – מתקין אותו מרצון (חושב שזה משחק/\enLR{crack)} & הוא עצמו – מנצל פגיעות ברשת אוטומטית & המשתמש (מריץ את הקובץ הנגוע, מעתיק בדיסק-און-קי) & \textbf{מי מפיץ?} \\
\hline
\textbf{לא} & כן, ובמהירות & כן & \textbf{משכפל את עצמו?} \\
\hline
\enLR{Zeus} (בנקאות), \enLR{RAT} (שליטה מרחוק) & \enLR{Morris, Slammer, WannaCry} & \enLR{Melissa (1999}, מאקרו ב-\enLR{Word)} & \textbf{דוגמה} \\
\hline
\end{xltabular}

\subheading{סוגי נוזקות נוספים שכדאי להכיר}

\begin{itemize}
\item \textbf{\enLR{Ransomware} (כופרה)} – מצפינה קבצים ודורשת תשלום. כיום הסוג המזיק ביותר לארגונים.
\item \textbf{\enLR{Spyware / Keylogger}} – מרגל, מתעד הקשות.
\item \textbf{\enLR{Adware}} – פרסומות; לרוב מטריד יותר ממסוכן.
\item \textbf{\enLR{Rootkit}} – מסתתר עמוק במערכת ההפעלה \enLR{(kernel)} ומסווה נוזקות אחרות; קשה מאוד לגילוי.
\item \textbf{\enLR{Botnet / Bot / Zombie}} – מחשבים נגועים הנשלטים מרחוק על ידי \enLR{C\&C (Command \& Control)} ומשמשים ל-\enLR{DDoS, spam}, כריית מטבעות.
\item \textbf{\enLR{Logic bomb}} – קוד ש"מתפוצץ" בתנאי מסוים (תאריך, פיטורי העובד).
\item \textbf{\enLR{Backdoor}} – דלת אחורית שעוקפת אימות. יכולה להיות מותקנת על ידי טרויאני או להישאר מכוונה בקוד.
\end{itemize}

\begin{faqbox}\boxtitle{שאלת תלמיד: "אנטי-וירוס מספיק?"}\par  אנטי-וירוס קלאסי מזהה לפי \textbf{חתימות} \enLR{(signature)} – תבניות של נוזקות ידועות. נוזקה חדשה \enLR{(zero-day)} או נוזקה שעברה שינוי קל \enLR{(polymorphic)} לא תזוהה. לכן כיום משלבים: ניתוח התנהגות \enLR{(heuristics), EDR}, עדכוני מערכת, הגנת רשת \enLR{(IPS)}, והכי חשוב – משתמש שלא לוחץ על כל דבר. "הגנה לעומק" \enLR{(Defense in Depth)} – כמה שכבות עצמאיות, כך שכשל באחת לא מפיל הכול.\end{faqbox}

\sectionheading{\enLR{1.5} חקר ההתקפות – שלושת השלבים}

מרבית התקפות הרשת מתחלקות לשלוש קטגוריות שסיסקו מגדירה, ואפשר גם לתאר אותן כשלבים בזמן:\par

\begin{defbox}\boxtitle{\enLR{1.} סיור / איסוף מידע \enLR{(Reconnaissance)} → \enLR{2.} השגת גישה \enLR{(Access)} → \enLR{3.} מניעת שירות \enLR{(DoS)} או ניצול}\par  מודל מפורט יותר, \textbf{\enLR{Cyber Kill Chain}} של לוקהיד-מרטין: \enLR{Reconnaissance} → \enLR{Weaponization} → \enLR{Delivery} → \enLR{Exploitation} → \enLR{Installation} → \enLR{Command \& Control} → \enLR{Actions on Objectives.} הרעיון: אם שוברים חוליה אחת בשרשרת – התקיפה נכשלת.\end{defbox}

\subheading{\enLR{1.5.1} איסוף מידע \enLR{(Reconnaissance)}}

התוקף לומד את המטרה לפני שנוגע בה. הבחינו בין:\par

\begin{itemize}
\item \textbf{סיור פסיבי} – ללא מגע עם המטרה: \enLR{whois} (מי הבעלים של הדומיין), \enLR{DNS} (\enLR{nslookup, dig}), רשתות חברתיות, אתר החברה, מודעות דרושים ("דרוש מנהל רשת עם ניסיון ב-\enLR{Cisco ASA"} – עכשיו התוקף יודע איזו חומת אש יש), \enLR{OSINT.}
\item \textbf{סיור אקטיבי} – שליחת חבילות:  \textbf{\enLR{Ping sweep}} – פינג לכל הכתובות בטווח כדי לגלות אילו מארחים חיים. \textbf{\enLR{Port scan}} – בדיקה אילו פורטים פתוחים בכל מארח (\enLR{Nmap}). פורט \enLR{22} פתוח ⇒ \enLR{SSH; 3389} ⇒ \enLR{RDP; 445} ⇒ \enLR{SMB.} \textbf{\enLR{Fingerprinting}} – זיהוי מערכת ההפעלה והגרסה \enLR{(Nmap -O, banner grabbing).} \textbf{\enLR{Packet sniffing}} – האזנה לתעבורה \enLR{(Wireshark)} – יעיל במיוחד ברשת אלחוטית פתוחה או כשמשתמשים בפרוטוקולים לא מוצפנים \enLR{(Telnet, HTTP, FTP).}
\end{itemize}

\textbf{הגנה:} \enLR{IPS} שמזהה סריקות, חומת אש שחוסמת \enLR{ICMP} מבחוץ, סגירת שירותים לא נחוצים, הצפנה נגד \enLR{sniffing}, מדיניות מידע ברשתות חברתיות.\par

\subheading{\enLR{1.5.2} התקפות גישה \enLR{(Access attacks)}}

\begin{xltabular}{\linewidth}{|>{\setRTL\raggedleft\arraybackslash}X|>{\setRTL\raggedleft\arraybackslash}X|>{\setRTL\raggedleft\arraybackslash}X|}
\hline
\rowcolor{hdrbg}\textbf{הגנה} & \textbf{איך זה עובד} & \textbf{התקפה} \\
\hline
\endhead
סיסמאות ארוכות ומורכבות, נעילה אחרי \enLR{X} ניסיונות, \enLR{MFA}, האטת ניסיונות & ניסוי שיטתי של כל הצירופים \enLR{(brute force)} או של רשימת מילים נפוצות \enLR{(dictionary).} כלים: \enLR{John the Ripper, Hydra.} בבגרות: "ניסיון לגלות סיסמה ע"י מילון" = \enLR{Brute force attack.} & \textbf{\enLR{Password attack}} – \enLR{Brute force / Dictionary} \\
\hline
מזעור יחסי אמון, פילוח רשת & ניצול יחסי אמון: אם שרת \enLR{A} סומך על \enLR{B}, פורצים ל-\enLR{B} (החלש) ומשם ל-\enLR{A} & \textbf{\enLR{Trust exploitation}} \\
\hline
\enLR{HIPS}, מודל אמון נכון בין \enLR{DMZ} לפנים & מארח שנפרץ (ב-\enLR{DMZ)} משמש כ"קופץ" להעברת תעבורה למארח פנימי שחומת האש לא מאפשרת אליו גישה ישירה & \textbf{\enLR{Port redirection}} \\
\hline
הצפנה עם אימות \enLR{(TLS, SSH} עם בדיקת \enLR{fingerprint), DAI} & התוקף מכניס את עצמו בין שני צדדים ומקשיב/משנה \enLR{(ARP spoofing} ברשת מקומית – פרק \enLR{6; Evil Twin} ב-\enLR{Wi-Fi)} & \textbf{\enLR{Man-in-the-Middle (MITM)}} \\
\hline
עדכוני תוכנה, \enLR{DEP/ASLR}, בדיקת קלט & שליחת קלט ארוך מהחוצץ שהוקצה לו, כך שהוא "גולש" ודורס זיכרון – כולל כתובת החזרה – ומריץ קוד של התוקף & \textbf{\enLR{Buffer overflow}} \\
\hline
הדרכה, נהלים & מניפולציה של בני אדם – ראו \enLR{1.6} & \textbf{\enLR{Social engineering}} \\
\hline
\end{xltabular}

\subheading{\enLR{1.5.3} מניעת שירות – \enLR{DoS} ו-\enLR{DDoS}}

\textbf{\enLR{DoS (Denial of Service)}:} הפיכת שירות ללא זמין למשתמשים הלגיטימיים – על ידי הצפה בתעבורה או ניצול באג שמקריס את השירות. \textbf{\enLR{DDoS (Distributed)}:} אותו דבר, אבל ממקורות רבים במקביל \enLR{(botnet)} – קשה לחסום כי אין כתובת מקור אחת.\par

\begin{xltabular}{\linewidth}{|>{\setRTL\raggedleft\arraybackslash}X|>{\setRTL\raggedleft\arraybackslash}X|}
\hline
\rowcolor{hdrbg}\textbf{מנגנון} & \textbf{התקפה} \\
\hline
\endhead
שליחת חבילת \enLR{ICMP} גדולה מ-\enLR{65,535} בתים (בעזרת פרגמנטציה); מערכות ישנות קרסו בהרכבה מחדש. היסטורי – היום מתוקן, אך נשאל בבחינות. & \textbf{\enLR{Ping of Death}} \\
\hline
שליחת המוני \enLR{SYN} בלי לסיים את לחיצת היד המשולשת; טבלת החיבורים "החצי-פתוחים" של השרת מתמלאת ולקוחות אמיתיים נדחים. הגנה: \enLR{SYN cookies, TCP intercept} בנתב. & \textbf{\enLR{SYN flood}} \\
\hline
שליחת פינג לכתובת \enLR{broadcast} עם כתובת מקור מזויפת של הקורבן; כל המחשבים ברשת עונים לקורבן (הגברה). הגנה: \code{no ip directed-broadcast} (ברירת מחדל היום). & \textbf{\enLR{Smurf}} \\
\hline
שאילתה קטנה עם \enLR{IP} מקור מזויף ⇒ תשובה גדולה נשלחת לקורבן. יחס הגברה של פי \enLR{50} ויותר. & \textbf{\enLR{Amplification / Reflection}} \enLR{(DNS, NTP)} \\
\hline
בקשות "לגיטימיות" רבות או חיבורים שנשארים פתוחים לאט – מתישים את השרת. & \textbf{\enLR{Application-layer DoS}} \enLR{(HTTP flood, Slowloris)} \\
\hline
\end{xltabular}

\textbf{הגנות מפני \enLR{DoS}:} סינון בכניסה לרשת \enLR{(anti-spoofing ACL} – פרק \enLR{4), rate limiting, IPS}, שירותי ניקוי תעבורה בענן \enLR{(Cloudflare, Akamai)}, יתירות ורוחב פס עודף, \code{ip verify unicast reverse-path} \enLR{(uRPF).}\par

\sectionheading{\enLR{1.6} הנדסה חברתית ו-\enLR{Phishing}}

ההתקפה שלא צריכה אף פגיעות טכנית – היא מנצלת את האדם. סוגים:\par

\begin{itemize}
\item \textbf{\enLR{Phishing}} – דוא"ל/\enLR{SMS} שמתחזה לגוף לגיטימי (בנק, \enLR{UPS}, משרד) ומבקש ללחוץ על קישור, "לאשר פרטים" או לפתוח קובץ. \textbf{\enLR{Spear phishing}} – ממוקד לאדם ספציפי עם פרטים אמיתיים עליו. \textbf{\enLR{Whaling}} – ממוקד למנכ"ל. \textbf{\enLR{Smishing / Vishing}} – ב-\enLR{SMS} / בטלפון.
\item \textbf{\enLR{Pretexting}} – סיפור כיסוי ("אני מהתמיכה, צריך את הסיסמה שלך כדי לתקן").
\item \textbf{\enLR{Baiting}} – דיסק-און-קי "אבוד" בחניה עם תווית "משכורות".
\item \textbf{\enLR{Tailgating / Piggybacking}} – כניסה פיזית בעקבות עובד מורשה.
\item \textbf{\enLR{Quid pro quo}} – "תן לי גישה ואתן לך מתנה".
\end{itemize}

\begin{exambox}\boxtitle{בבגרות (שאלה \enLR{4}א): איך מזהים הודעת \enLR{phishing}? סימנים לחפש}\par  \enLR{1.} כתובת השולח לא תואמת לארגון \enLR{(UPS} ששולח מ-\enLR{@hotmail.com}). \enLR{2.} תחושת דחיפות/איום ("\enLR{Last Call}!", "החשבון ייחסם"). \enLR{3.} בקשה לאשר/להזין פרטים דרך כפתור. \enLR{4.} שגיאות ניסוח, לוגו מטושטש, פנייה כללית ("\enLR{Dear customer"). 5.} קישור שהטקסט שלו שונה מהיעד האמיתי (מרחפים עם העכבר). \textbf{התשובה: \enLR{Social Engineering}}.\end{exambox}

\begin{storybox}\boxtitle{סיפור מהחיים: הסמנכ"ל ששילם \enLR{40} מיליון (הונאת "מנכ"ל", \enLR{2016)}}\par  עובד כספים בחברת \enLR{Crelan} הבלגית קיבל מייל "מהמנכ"ל": "אנחנו ברכישה סודית, העבר דחוף \enLR{70} מיליון אירו לחשבון הזה, אל תדבר עם אף אחד". הכתובת הייתה דומה לכתובת המנכ"ל בהבדל אות אחת. הכסף הועבר. זו הנדסה חברתית מסוג \textbf{\enLR{BEC}} \enLR{(Business Email Compromise)} – בלי נוזקה, בלי פריצה, רק פסיכולוגיה: סמכות + דחיפות + סודיות. שאלו את התלמידים: איזה נוהל (בקרה מנהלית) היה מונע את זה? (אישור כפול לכל העברה מעל סכום, ואימות טלפוני בערוץ נפרד.)\end{storybox}

\sectionheading{\enLR{1.7} שיכוך מתקפות ברוט-פורס – ומה אפשר לעשות כבר בנתב}

"ברוטלי" בתכנית = \enLR{Brute force.} עקרונות ההגנה:\par

\begin{enumerate}
\item \textbf{סיסמאות חזקות} – אורך חשוב יותר ממורכבות. \enLR{8} תווים מאותיות קטנות = \enLR{26}\textsuperscript{\enLR{8}} ≈ \enLR{200} מיליארד (נפרץ בדקות ב-\enLR{GPU); 12} תווים מכל הסוגים = \enLR{95}\textsuperscript{\enLR{12}} ≈ \enLR{5}×\enLR{10}\textsuperscript{\enLR{23}}. משפט-סיסמה \enLR{(passphrase)} עדיף.
\item \textbf{נעילה והאטה} – אחרי \enLR{N} ניסיונות כושלים: נעילת חשבון, השהיה מתארכת \enLR{(exponential backoff), CAPTCHA.}
\item \textbf{אימות רב-שלבי \enLR{(MFA)}} – גם אם הסיסמה נפרצה, התוקף לא מחזיק את הטלפון.
\item \textbf{אחסון סיסמאות נכון} – לא בטקסט גלוי; גיבוב עם \enLR{salt} (פרק \enLR{7).}
\item \textbf{ניטור} – התראה על ריבוי כניסות כושלות \enLR{(syslog} – פרק \enLR{2).}
\end{enumerate}

ב-\enLR{Cisco IOS} יש מנגנון מובנה שנלמד לעומק בפרק \enLR{2}, אבל שווה להציג אותו כבר כאן כדי לחבר תיאוריה למעשה:\par

\begin{codebox}\begin{english}\codefont\footnotesize\color{cfg}\setlength{\parindent}{0pt}
R1(config)\# \textcolor{cbold}{\textbf{login block-for 120 attempts 3 within 60}}\\
\textcolor{ccom}{\texthebrew{! אם היו \enLR{3} ניסיונות כושלים תוך \enLR{60} שניות – חסום כניסות ל-\enLR{120} שניות}}\\
R1(config)\# \textcolor{cbold}{\textbf{login delay 2}}          \textcolor{ccom}{\texthebrew{! השהיה של \enLR{2} שניות בין ניסיונות}}\\
R1(config)\# \textcolor{cbold}{\textbf{login on-failure log}}    \textcolor{ccom}{\texthebrew{! רשום כל כישלון ל-\enLR{syslog}}}\\
R1(config)\# \textcolor{cbold}{\textbf{login on-success log}}\\
R1\# \textcolor{cbold}{\textbf{show login}}
\end{english}\end{codebox}

\sectionheading{\enLR{1.8} ארגון האבטחה: מדיניות, תפקידים ותקנים}

\subheading{מדיניות אבטחה \enLR{(Security Policy)}}

מסמך שקובע \emph{מה} מותר ומה אסור בארגון – ולא \emph{איך} (זה בנהלים). מרכיבים אופייניים: מדיניות שימוש מקובל \enLR{(AUP)}, מדיניות סיסמאות, גישה מרחוק, סיווג מידע, תגובה לאירועים. בלי מדיניות כתובה אי אפשר לאכוף ולא ניתן לפטר עובד ש"רק בדק".\par

\subheading{סוגי בקרות (נשאל בבגרות – שאלה \enLR{6}ב)}

\begin{xltabular}{\linewidth}{|>{\setRTL\raggedleft\arraybackslash}X|>{\setRTL\raggedleft\arraybackslash}X|>{\setRTL\raggedleft\arraybackslash}X|}
\hline
\rowcolor{hdrbg}\textbf{דוגמאות} & \textbf{תיאור} & \textbf{סוג בקרה} \\
\hline
\endhead
איסור על עבודה ב-\enLR{Zoom} מרשת לא מאובטחת, נוהל סיסמאות, הכשרת עובדים & מדיניות, נהלים, הנחיות, תקנים, הדרכות, בדיקות רקע & \textbf{מנהלית} \enLR{(Administrative)} \\
\hline
מנעולים על חדר שרתים וארונות תקשורת, מצלמות, שומרים, כרטיסי כניסה & הגנה על החומרה והמבנה & \textbf{פיזית} \enLR{(Physical)} \\
\hline
רשימות גישה \enLR{(ACL)}, סיסמאות, הצפנה, חומת אש, \enLR{IPS}, אנטי-וירוס & אמצעים בתוכנה ובחומרה & \textbf{לוגית / טכנית} \enLR{(Logical / Technical)} \\
\hline
\end{xltabular}

\subheading{תפקידים}

\textbf{\enLR{CISO}} \enLR{(Chief Information Security Officer)} – אחראי אבטחה ארגוני; \textbf{\enLR{SOC}} \enLR{(Security Operations Center)} – מרכז ניטור \enLR{24/7}; \textbf{\enLR{CERT/CSIRT}} – צוות תגובה לאירועים; בישראל: \textbf{מערך הסייבר הלאומי} ו-\enLR{CERT-IL.}\par

\subheading{תקנים וארגונים}

\begin{itemize}
\item \textbf{\enLR{ISO/IEC 27001}} – תקן בינלאומי לניהול אבטחת מידע \enLR{(ISMS)}; ארגונים מוסמכים לו.
\item \textbf{\enLR{NIST Cybersecurity Framework}} – \enLR{Identify, Protect, Detect, Respond, Recover.}
\item \textbf{\enLR{PCI-DSS}} – חובה לכל מי שמעבד כרטיסי אשראי.
\item \textbf{\enLR{CVE}} \enLR{(Common Vulnerabilities and Exposures)} – מאגר עולמי של פגיעויות עם מזהה, לדוגמה \enLR{CVE-2017-0144} (הפגיעות ש-\enLR{WannaCry} ניצלה). \textbf{\enLR{CVSS}} – ציון חומרה \enLR{0}–\enLR{10.}
\item \textbf{\enLR{SANS, OWASP}} \enLR{(Top 10} לאפליקציות \enLR{web)}, \textbf{\enLR{MITRE ATT\&CK}} – קטלוג טקטיקות תוקפים.
\end{itemize}

\sectionheading{\enLR{1.9} אביזרי אבטחה ברשת – מפת הקורס}

סקירה של הרכיבים שנלמד בפרקים הבאים, כדי שהתלמידים יראו את התמונה הגדולה:\par

\begin{xltabular}{\linewidth}{|>{\setRTL\raggedleft\arraybackslash}X|>{\setRTL\raggedleft\arraybackslash}X|>{\setRTL\raggedleft\arraybackslash}X|}
\hline
\rowcolor{hdrbg}\textbf{פרק} & \textbf{תפקיד} & \textbf{רכיב} \\
\hline
\endhead
\enLR{2} & קו ההגנה הראשון בקצה; \enLR{SSH}, סיסמאות, \enLR{logging} & נתב מוקשח \enLR{(Hardened router)} \\
\hline
\enLR{3} & אימות מרכזי של משתמשים ומנהלים & שרת \enLR{AAA (RADIUS/TACACS+)} \\
\hline
\enLR{4} & סינון תעבורה לפי כללים & חומת אש \enLR{(Firewall) / ACL / ZPF} \\
\hline
\enLR{5} & זיהוי וחסימת התקפות לפי חתימות & \enLR{IDS / IPS} \\
\hline
\enLR{6} & הגנה בשכבה \enLR{2} & מתג מאובטח \enLR{(port security, DHCP snooping, DAI)} \\
\hline
\enLR{7} & סודיות ושלמות & הצפנה, \enLR{PKI} \\
\hline
\enLR{8} & ערוץ מוצפן על גבי רשת ציבורית & \enLR{VPN (IPsec, SSL)} \\
\hline
\enLR{9} & בדיקה שהכול עובד & סורקי פגיעויות, ניהול סיכונים \\
\hline
\end{xltabular}

\textbf{\enLR{Cisco ASA}} \enLR{(Adaptive Security Appliance)} – מכשיר ייעודי שמשלב חומת אש, \enLR{VPN} ו-\enLR{IPS.} יורשו הנוכחי: \textbf{\enLR{Firepower / FTD}}. בתכנית מוזכר גם \textbf{\enLR{Cisco SDM}} – ממשק ניהול גרפי ישן; היום \enLR{CLI.}\par

\begin{mistakebox}\boxtitle{טעויות נפוצות של תלמידים בפרק זה}\par  • מבלבלים וירוס ותולעת – השאלה המכריעה: \textbf{האם צריך פעולה של משתמש כדי להתפשט?} וירוס – כן. תולעת – לא.\newline  • חושבים שטרויאני משכפל את עצמו – לא. הוא מחכה שיתקינו אותו.\newline  • "\enLR{DDoS} זה וירוס" – \enLR{DDoS} היא \emph{התקפה}; \enLR{botnet} (שנבנה בעזרת נוזקה) הוא \emph{הכלי}.\newline  • "\enLR{Phishing} זה וירוס במייל" – \enLR{phishing} הוא הנדסה חברתית; המייל עצמו לא בהכרח מכיל נוזקה, לפעמים רק קישור לאתר מזויף.\newline  • "\enLR{Hacker} = פושע" – הבהירו את \enLR{white/black/gray hat.}\end{mistakebox}

\sectionheading{\enLR{1.10} תרגילים לתלמידים}

\begin{exbox}\boxtitle{תרגיל \enLR{1} – סיווג אירועים לפי \enLR{CIA}}\par  לכל אירוע, ציינו איזה עיקרון נפגע (יכול להיות יותר מאחד):\newline  א. תלמיד מצא דרך לשנות ציון במערכת בית הספר. ב. אתר בית הספר קרס ביום ההרשמה. ג. רשימת הטלפונים של ההורים דלפה לקבוצת וואטסאפ. ד. כופרה הצפינה את שרת הקבצים. ה. מישהו האזין לשיחת \enLR{Zoom} שלא הוזמן אליה. \par\vskip2pt\hrule\vskip3pt\textbf{}א. שלמות ב. זמינות ג. סודיות ד. זמינות (ואם הם גם מאיימים לפרסם – סודיות) ה. סודיות\end{exbox}

\begin{exbox}\boxtitle{תרגיל \enLR{2} – וירוס, תולעת או טרויאני?}\par  א. תוכנה בשם "\enLR{Free-Minecraft-Coins.exe"} שמתקינה \enLR{keylogger.} ב. קובץ \enLR{Excel} עם מאקרו שמדביק את כל קובצי ה-\enLR{Excel} האחרים במחשב כשפותחים אותו. ג. תוכנית שסורקת את הרשת, מוצאת מחשבים עם פורט \enLR{445} פתוח ומעתיקה את עצמה אליהם. ד. אפליקציית פנס לאנדרואיד ששולחת את אנשי הקשר לשרת בסין. \par\vskip2pt\hrule\vskip3pt\textbf{}א. טרויאני ב. וירוס (מאקרו) ג. תולעת ד. טרויאני \enLR{(spyware)}\end{exbox}

\begin{exbox}\boxtitle{תרגיל \enLR{3} – חישוב חוזק סיסמה}\par  מחשב תקיפה מנסה \enLR{10} מיליארד \enLR{(10}\textsuperscript{\enLR{10}}) סיסמאות בשנייה. כמה זמן ייקח במקרה הגרוע לפרוץ: א. \enLR{6} ספרות \enLR{(PIN).} ב. \enLR{8} אותיות קטנות. ג. \enLR{8} תווים מכל \enLR{95} התווים. ד. \enLR{12} תווים מכל \enLR{95} התווים. (חישוב: מספר האפשרויות ÷ \enLR{10}\textsuperscript{\enLR{10}}) \par\vskip2pt\hrule\vskip3pt\textbf{}א. \enLR{10}\textsuperscript{\enLR{6}}/\enLR{10}\textsuperscript{\enLR{10}} = \enLR{0.0001} שנייה. ב. \enLR{26}\textsuperscript{\enLR{8}} ≈ \enLR{2.1}×\enLR{10}\textsuperscript{\enLR{11}} ⇒ \textasciitilde{}\enLR{21} שניות. ג. \enLR{95}\textsuperscript{\enLR{8}} ≈ \enLR{6.6}×\enLR{10}\textsuperscript{\enLR{15}} ⇒ \textasciitilde{}\enLR{7.6} ימים. ד. \enLR{95}\textsuperscript{\enLR{12}} ≈ \enLR{5.4}×\enLR{10}\textsuperscript{\enLR{23}} ⇒ \textasciitilde{}\enLR{1.7} מיליון שנים. מסקנה: כל תו נוסף מכפיל את הזמן פי \enLR{95.}\end{exbox}

\begin{exbox}\boxtitle{תרגיל \enLR{4} – זהו את שלב התקיפה}\par  לתוקף יש את הרצף הבא. סדרו לפי \enLR{Reconnaissance / Access / DoS}, וציינו לכל שלב מה היה אפשר לעשות כדי לעצור אותו:\newline  \enLR{1.} הריץ \enLR{Nmap} על טווח הכתובות של החברה. \enLR{2.} מצא שרת \enLR{FTP} ישן וניחש את הסיסמה \enLR{admin/admin. 3.} העלה לשרת סקריפט שמציף את הראוטר בבקשות. \enLR{4.} קרא במודעת דרושים שהחברה משתמשת ב-\enLR{Windows Server 2012.} \par\vskip2pt\hrule\vskip3pt\textbf{}\enLR{Reconnaissance: 4} (פסיבי), \enLR{1} (אקטיבי) – הגנה: מדיניות פרסום, \enLR{IPS} שמזהה סריקות. \enLR{Access: 2} – הגנה: סיסמאות חזקות, כיבוי שירותים מיושנים, \enLR{login block. DoS: 3} – הגנה: \enLR{rate limiting, IPS, HIPS} על השרת.\end{exbox}

\begin{exbox}\boxtitle{תרגיל \enLR{5} – דיון: "אני לא מעניין אף אחד"}\par  תלמיד טוען: "אין לי כלום במחשב, למה שמישהו יתקוף אותי?" תנו \enLR{4} סיבות שבגללן גם מחשב "לא מעניין" הוא מטרה. \par\vskip2pt\hrule\vskip3pt\textbf{}\enLR{1.} הוא יכול להפוך ל-\enLR{bot} ב-\enLR{botnet (DDoS, spam}, כריית מטבעות). \enLR{2.} הסיסמאות שלו כנראה זהות לסיסמאות בחשבונות אחרים \enLR{(credential stuffing). 3.} הוא קפיצת מדרגה לרשת של ההורים/בית הספר. \enLR{4.} הזהות שלו (חשבון אינסטגרם, מייל) שווה כסף להונאות. \enLR{5.} כופרה לא בוררת – היא אוטומטית.\end{exbox}

\sectionheading{\enLR{1.11} שאלות בסגנון בגרות – עם פתרונות}

\begin{exambox}\boxtitle{שאלה \enLR{1}}\par  משתמש קיבל \enLR{SMS: "}חבילתך ממתינה. לתשלום מכס של \enLR{4.90} ₪ לחץ כאן: \enLR{post-il.delivery-confirm.xyz".} מהו סוג ההתקפה?\newline \enLR{1. Ping of Death 2. Smishing} (הנדסה חברתית) \enLR{3. Buffer overflow 4.} תולעת \par\vskip2pt\hrule\vskip3pt\textbf{}\textbf{\enLR{2.}} הודעת \enLR{SMS} מתחזה עם דחיפות וסכום קטן שנועד להשיג פרטי כרטיס אשראי. הדומיין אינו של דואר ישראל.\end{exambox}

\begin{exambox}\boxtitle{שאלה \enLR{2}}\par  התאימו: (א) קוד שמצמיד את עצמו לקובץ ומתפשט כשמריצים אותו; (ב) תוכנית שמתפשטת ברשת ללא התערבות משתמש; (ג) תוכנה שנראית כמשחק אך פותחת דלת אחורית. אפשרויות: \enLR{Worm / Trojan / Virus.} \par\vskip2pt\hrule\vskip3pt\textbf{}(א) \enLR{Virus} (ב) \enLR{Worm} (ג) \enLR{Trojan}\end{exambox}

\begin{exambox}\boxtitle{שאלה \enLR{3}}\par  תוקף שולח אלפי חבילות \enLR{SYN} לשרת ואינו משיב ל-\enLR{SYN-ACK.} איזו התקפה זו ואיזה עיקרון ב-\enLR{CIA} נפגע? \par\vskip2pt\hrule\vskip3pt\textbf{}\enLR{SYN flood (DoS).} נפגעת ה\textbf{זמינות} \enLR{(Availability).} טבלת החיבורים החצי-פתוחים מתמלאת.\end{exambox}

\begin{exambox}\boxtitle{שאלה \enLR{4}}\par  מנהל רשת רוצה שהנתב יחסום ניסיונות התחברות ל-\enLR{60} שניות לאחר \enLR{4} ניסיונות כושלים בתוך \enLR{30} שניות. כתבו את הפקודה. \par\vskip2pt\hrule\vskip3pt\textbf{}\code{login block-for 60 attempts 4 within 30}\end{exambox}

\begin{exambox}\boxtitle{שאלה \enLR{5}}\par  עובד לשעבר שעדיין מחזיק בסיסמה ל-\enLR{VPN} של החברה ומוריד ממנה קבצים – לאיזה סוג תוקף הוא שייך? מה הבקרה (מנהלית) שהייתה מונעת זאת? \par\vskip2pt\hrule\vskip3pt\textbf{}\enLR{Insider.} נוהל סיום העסקה \enLR{(offboarding)} שכולל ביטול מיידי של כל החשבונות – בקרה מנהלית; ביטול החשבון בפועל – בקרה לוגית.\end{exambox}

\begin{exambox}\boxtitle{שאלה \enLR{6}}\par  איזה שלב בתקיפה כולל \enLR{ping sweep} ו-\enLR{port scan}? מה הכלי המקובל? \par\vskip2pt\hrule\vskip3pt\textbf{}\enLR{Reconnaissance} (סיור/איסוף מידע). \enLR{Nmap.}\end{exambox}

\sectionheading{\enLR{1.12} הצעה למעבדה / פעילות}

\begin{itemize}
\item \textbf{\enLR{Wireshark} – \enLR{Telnet} מול \enLR{SSH}:} הריצו \enLR{Telnet} לנתב ב-\enLR{Packet Tracer} (או במעבדה) ולכדו את התעבורה; הראו את הסיסמה בטקסט גלוי. חזרו על זה עם \enLR{SSH.} זה מכין את פרק \enLR{2.}
\item \textbf{ניתוח \enLR{5} הודעות דוא"ל:} תנו לתלמידים \enLR{5} צילומי מסך \enLR{(3 phishing, 2} לגיטימיות) ובקשו לסמן את הסימנים. דוגמאות אמיתיות (מטושטשות) עובדות הכי טוב.
\item \textbf{חישוב חוזק סיסמה:} חשבו עם התלמידים כמה זמן ייקח לפרוץ סיסמה ב-\enLR{10}\textsuperscript{\enLR{10}} ניסיונות/שנייה עבור אורכים שונים. המסקנה – אורך מנצח.
\end{itemize}

\sectionheading{עוד דוגמאות מהעולם האמיתי}

\begin{storybox}\boxtitle{\enLR{NotPetya (2017)} – הכופרה שלא הייתה כופרה, נזק של \enLR{10} מיליארד דולר}\par  \enLR{NotPetya} נראתה כמו כופרה (דרשה תשלום), אך למעשה הייתה \textbf{\enLR{Wiper}} – מחקה את המידע לצמיתות בלי אפשרות שחזור. היא התפשטה דרך \underline{עדכון תוכנה לגיטימי} של תוכנת הנהלת חשבונות אוקראינית פופולרית \enLR{(M.E.Doc)} – מתקפת \textbf{שרשרת אספקה}. תוך שעות שיתקה חברות ענק: ענקית הספנות \enLR{Maersk} (שנאלצה להתקין מחדש \enLR{4,000} שרתים ו-\enLR{45,000} מחשבים), חברת התרופות \enLR{Merck}, ו-\enLR{FedEx.} הנזק הכולל נאמד ב-\enLR{10} מיליארד דולר. \textbf{לקח לכיתה:} נוזקה יכולה להיכנס דרך תוכנה שאתם סומכים עליה ומתעדכנת אוטומטית – לא רק דרך קליק שלכם. וגם: לא כל מה שנראה ככופרה נועד לכסף; לפעמים המטרה היא הרס טהור (מתקפה מדינתית).\end{storybox}

\begin{storybox}\boxtitle{פריצת טוויטר של המפורסמים \enLR{(2020)} – הנדסה חברתית בטלפון}\par  קבוצת בני נוער התקשרה לעובדי טוויטר, התחזתה לצוות ה-\enLR{IT} הפנימי, ו\underline{שכנעה אותם למסור פרטי גישה} לכלי ניהול פנימי \enLR{(vishing} – \enLR{phishing} טלפוני). עם הגישה הזו הם השתלטו על חשבונות של ברק אובמה, אילון מאסק, ביל גייטס, אפל ו-\enLR{Uber}, ופרסמו מהם הונאת ביטקוין ("שלחו ביטקוין ואחזיר כפול"). הם גרפו כ-\enLR{118,000}\$ תוך שעות – \textbf{בלי לפרוץ אף מחשב}, רק בכוח השכנוע. לקח: אפילו אחת מחברות הטכנולוגיה הגדולות בעולם נפלה להנדסה חברתית. האדם הוא החוליה החלשה, תמיד.\end{storybox}

\begin{storybox}\boxtitle{\enLR{Equifax (2017)} – פגיעות אחת שלא תוקנה, \enLR{147} מיליון אנשים}\par  חברת דירוג האשראי \enLR{Equifax} איבדה מידע רגיש (מספרי ביטוח לאומי, תאריכי לידה) של \enLR{147} מיליון אמריקאים. הסיבה: פגיעות ידועה בספריית \enLR{Apache Struts}, עם \underline{תיקון שהיה זמין חודשיים קודם} ולא הותקן. סורק פגיעויות (פרק \enLR{9)} היה מוצא אותה בדקות. לקח: ניהול עדכונים אינו אירוע חד-פעמי אלא תהליך מתמשך – בדיוק כפי שראינו ב-\enLR{WannaCry.}\end{storybox}

\begin{exbox}\boxtitle{תרגיל \enLR{6} – זהו את סוג המתקפה}\par  לכל תרחיש, אילו מהמושגים מהפרק מתאימים? (יכול להיות יותר מאחד)\newline  א. עדכון תוכנה לגיטימי הפיץ נוזקה מחקת מידע. ב. עובד התקשר והתחזה לתמיכה כדי לקבל סיסמה. ג. חצי מיליון מכשירי \enLR{IoT} הפכו לצבא שהפיל אתרים. \par\vskip2pt\hrule\vskip3pt\textbf{}א. מתקפת שרשרת אספקה + \enLR{Wiper} (מדינתי) – \enLR{NotPetya.} ב. הנדסה חברתית / \enLR{vishing} – פריצת טוויטר. ג. \enLR{botnet + DDoS} – \enLR{Mirai.}\end{exbox}

\begin{exbox}\boxtitle{תרגיל \enLR{7} – חשיבה}\par  בכל שלושת הסיפורים לעיל, איזו בקרה (מנהלית / פיזית / לוגית) הייתה מקטינה את הנזק? תנו דוגמה אחת לכל סיפור. \par\vskip2pt\hrule\vskip3pt\textbf{}\enLR{NotPetya}: בקרה לוגית – פילוח רשת וגיבוי מנותק \enLR{(3-2-1)} לשחזור. טוויטר: בקרה מנהלית – נוהל שאוסר מסירת גישה בטלפון + אימות זהות; ולוגית – \enLR{MFA} וכלי ניהול עם הרשאות מוגבלות. \enLR{Equifax}: מנהלית – נוהל עדכונים ובדיקות פגיעות תקופתיות.\end{exbox}

\sectionheading{בנק שאלות שתלמידים שואלים – ותשובות מוכנות}

אוסף שאלות שתלמידים סקרנים נוטים לשאול בפרק זה, עם תשובות מוכנות. מסודר לפי נושא – סרקו לפני השיעור.\par

\subheading{מוטיבציה ותפיסות שגויות}

\textbf{ש: "אין לי מה להסתיר / אני לא מעניין אף אחד, אז למה שיתקפו אותי?"}\newline ת: רוב ההתקפות אינן ממוקדות אלא \underline{אוטומטיות וסיטונאיות}. גם מחשב "לא מעניין" שווה לתוקף: אפשר להפוך אותו ל-\enLR{bot} ברשת זומבים (פרק \enLR{1.4)}, לגנוב סיסמאות שחוזרות על עצמן באתרים אחרים \enLR{(credential stuffing)}, להשתמש בו כקרש קפיצה לרשת של בית הספר/ההורים, או להצפין אותו בכופרה ולדרוש כסף. כופרה לא בוררת מטרות.\par

\textbf{ש: "אנטי-וירוס מספיק כדי להיות מוגן?"}\newline ת: לא. אנטי-וירוס קלאסי מזהה לפי \textbf{חתימות} – תבניות של נוזקות \underline{ידועות}. נוזקה חדשה \enLR{(zero-day)} או כזו ששונתה מעט \enLR{(polymorphic)} לא תזוהה. לכן צריך "הגנה לעומק": עדכוני מערכת, הגנת רשת (חומת אש/\enLR{IPS)}, גיבוי, והכי חשוב – משתמש שלא לוחץ על כל דבר.\par

\textbf{ש: "מה זה \enLR{zero-day}?"}\newline ת: פגיעות שהיצרן עדיין \underline{לא יודע עליה} או שאין לה תיקון – "יום אפס" מאז שהתגלתה. תוקף שמחזיק \enLR{zero-day} יכול לנצל אותה בלי שאף הגנה מבוססת-חתימות תזהה אותה. אלה הפגיעויות היקרות והמסוכנות ביותר.\par

\textbf{ש: "האקר = פושע?"}\newline ת: לא בהכרח. \textbf{\enLR{White hat}} – בודק חדירה חוקי, באישור. \textbf{\enLR{Black hat}} – פושע. \textbf{\enLR{Gray hat}} – פורץ ללא אישור אך מדווח. המילה "\enLR{hacker"} במקורה היא מי שמבין מערכות לעומק; התקשורת הצמידה לה משמעות שלילית.\par

\subheading{נוזקות}

\textbf{ש: "אז מה בדיוק ההבדל בין וירוס לתולעת? זה מבלבל."}\newline ת: שאלת המפתח: \underline{האם צריך פעולה של משתמש כדי שיתפשט?} וירוס נדבק לקובץ וזקוק למשתמש שיריץ אותו. תולעת מתפשטת \underline{לבד} דרך הרשת בלי אף פעולת אדם (למשל תולעת מוריס, \enLR{Slammer).} טרויאני לא משכפל את עצמו כלל – המשתמש מתקין אותו מרצון כי הוא נראה לגיטימי.\par

\textbf{ש: "האם מק / לינוקס / אייפון יכולים להידבק בנוזקה?"}\newline ת: כן, כל מערכת יכולה. \enLR{Windows} הוא היעד הנפוץ ביותר כי הוא הכי נפוץ, אבל יש נוזקות ל-\enLR{macOS}, לינוקס, אנדרואיד ו-\enLR{iOS.} אף מערכת אינה "חסינה" – זו תפיסה שגויה מסוכנת.\par

\textbf{ש: "כדאי לשלם כופר בכופרה?"}\newline ת: הגישה הרשמית (גם של ה-\enLR{FBI} ומערך הסייבר) היא \underline{לא לשלם}: אין ערובה שתקבלו מפתח, זה מממן פשע, והופך אתכם ליעד חוזר. הפתרון האמיתי הוא \textbf{גיבוי} תקין (פרק \enLR{9)} שמאפשר לשחזר בלי לשלם.\par

\subheading{התקפות והגנה}

\textbf{ש: "אם תוקף מפיל אותנו ב-\enLR{DDoS}, למה שלא פשוט נחסום את הכתובת שלו?"}\newline ת: כי ב-\enLR{DDoS} אין כתובת אחת – ההתקפה מגיעה בו-זמנית מאלפי מכשירים נגועים \enLR{(botnet)} בכל העולם, ולעיתים עם כתובות מקור \underline{מזויפות}. חסימה ידנית לא מספיקה; צריך סינון בקנה מידה גדול, שירותי ניקוי בענן, ו-\enLR{rate limiting.}\par

\textbf{ש: "אדם אחד יכול לבצע \enLR{DDoS}?"}\newline ת: כן – אם הוא שולט ב-\enLR{botnet} (רשת מחשבים נגועים) או שוכר אחד בדארקנט ("\enLR{DDoS} כשירות"). ההבחנה היא לא מספר האנשים אלא מספר \underline{מקורות} התעבורה: \enLR{DoS} = מקור אחד, \enLR{DDoS} = מקורות רבים.\par

\textbf{ש: "איך \enLR{phishing} עוקף אנטי-וירוס?"}\newline ת: כי לרוב אין בו נוזקה לסרוק – יש רק קישור לאתר מזויף או בקשה להזין פרטים. זו מניפולציה של \underline{האדם}, לא של המחשב. אנטי-וירוס לא "רואה" שאתם מקלידים סיסמה באתר מתחזה.\par

\textbf{ש: "אם יש נעילה אחרי \enLR{3} ניסיונות, למה צריך בכלל סיסמה חזקה?"}\newline ת: כי לא תמיד התוקף עובר דרך מסך ההתחברות. אם דלף \underline{מסד הנתונים} של הסיסמאות (גיבובים), התוקף מריץ מיליארדי ניחושים אופליין על החומרה שלו – שם אין נעילה. סיסמה ארוכה הופכת את זה לבלתי מעשי.\par

\textbf{ש: "האם זה חוקי להריץ \enLR{Nmap} או להתקין \enLR{Kali Linux}?"}\newline ת: הכלים עצמם חוקיים לחלוטין (זה כמו פטיש). מה שאסור הוא \underline{השימוש} בהם על רשת שאינה שלכם וללא אישור בכתב – זו עבירה על חוק המחשבים התשנ"ה-\enLR{1995.} ברשת מעבדה מבודדת שלכם – מותר ומומלץ ללמוד.\par

\textbf{ש: "מה ההבדל המדויק בין איום, פגיעות וסיכון?"}\newline ת: \textbf{איום} = הגורם המסוכן (האקר, שריפה) – קיים תמיד, לא בשליטתנו. \textbf{פגיעות} = החולשה שהאיום מנצל (סיסמה חלשה, פורט פתוח) – זה מה שבשליטתנו. \textbf{סיכון} = ההסתברות שהאיום ינצל את הפגיעות, כפול הנזק. מנהלים סיכון על ידי צמצום פגיעויות.\par

\sectionheading{מילון מונחים – הגדרות}

הגדרות תמציתיות של המונחים המרכזיים בפרק. שימושי גם כדף עזר לבחינה (מותר מילון מונחים).\par

\begin{xltabular}{\linewidth}{|>{\setRTL\raggedleft\arraybackslash}X|>{\setRTL\raggedleft\arraybackslash}X|}
\hline
\rowcolor{hdrbg}\textbf{הגדרה} & \textbf{מונח} \\
\hline
\endhead
הגנה על מידע ומערכות מפני גישה, שינוי, חשיפה או השמדה בלתי מורשים. & \textbf{אבטחת מידע \enLR{(Information Security)}} \\
\hline
תת-תחום: הגנה על תשתית הרשת ועל התעבורה שעוברת בה. & \textbf{אבטחת רשת \enLR{(Network Security)}} \\
\hline
הבטחה שרק גורמים מורשים רואים את המידע. מושגת בהצפנה ובבקרת גישה. & \textbf{סודיות \enLR{(Confidentiality)}} \\
\hline
הבטחה שהמידע לא שונה ללא הרשאה. מושגת בגיבוב וחתימה דיגיטלית. & \textbf{שלמות \enLR{(Integrity)}} \\
\hline
הבטחה שהמידע והשירות נגישים כשצריך. נפגעת ב-\enLR{DoS} ובכופרה. & \textbf{זמינות \enLR{(Availability)}} \\
\hline
הוכחת זהות – מי אתה (שם משתמש וסיסמה, ביומטריה). & \textbf{אימות \enLR{(Authentication)}} \\
\hline
מניעת אפשרות של השולח להכחיש ששלח; מושגת בחתימה דיגיטלית. & \textbf{אי-הכחשה \enLR{(Non-repudiation)}} \\
\hline
כל דבר בעל ערך שמגנים עליו: מידע, שרת, מוניטין, זמינות. & \textbf{נכס \enLR{(Asset)}} \\
\hline
חולשה שניתן לנצל: באג, סיסמה חלשה, פורט פתוח, עובד לא מודרך. & \textbf{פגיעות \enLR{(Vulnerability)}} \\
\hline
גורם שעלול לנצל פגיעות: האקר, נוזקה, עובד ממורמר, אסון טבע. & \textbf{איום \enLR{(Threat)}} \\
\hline
הסתברות שאיום ינצל פגיעות, כפול הנזק. = איום × פגיעות × ערך הנכס. & \textbf{סיכון \enLR{(Risk)}} \\
\hline
אמצעי שמקטין סיכון. מנהלית (נהלים), פיזית (מנעולים), לוגית (הצפנה, \enLR{ACL).} & \textbf{בקרה \enLR{(Control)}} \\
\hline
שם כולל לתוכנה זדונית: וירוס, תולעת, טרויאני, כופרה ועוד. & \textbf{נוזקה \enLR{(Malware)}} \\
\hline
קוד זדוני שנצמד לקובץ ומתפשט כשמשתמש מריץ אותו. & \textbf{וירוס \enLR{(Virus)}} \\
\hline
תוכנה שמשכפלת ומתפשטת לבד ברשת, ללא פעולת משתמש. & \textbf{תולעת \enLR{(Worm)}} \\
\hline
תוכנה שנראית לגיטימית אך מכילה קוד זדוני; אינה משכפלת את עצמה. & \textbf{סוס טרויאני \enLR{(Trojan)}} \\
\hline
נוזקה שמצפינה קבצים ודורשת תשלום; פוגעת בזמינות. & \textbf{כופרה \enLR{(Ransomware)}} \\
\hline
נוזקה שמסתתרת עמוק במערכת ההפעלה ומסווה נוזקות אחרות. & \textbf{\enLR{Rootkit}} \\
\hline
רשת מחשבים נגועים (זומבים) הנשלטים מרחוק לביצוע התקפות. & \textbf{\enLR{Botnet}} \\
\hline
דלת אחורית שעוקפת את מנגנון האימות ומאפשרת גישה סמויה. & \textbf{\enLR{Backdoor}} \\
\hline
מניפולציה של בני אדם כדי להשיג מידע או גישה (במקום פריצה טכנית). & \textbf{הנדסה חברתית \enLR{(Social Engineering)}} \\
\hline
דוא"ל/הודעה מתחזים המבקשים ללחוץ על קישור או למסור פרטים. & \textbf{\enLR{Phishing}} \\
\hline
מניעת שירות – הצפת שירות עד קריסה, ממקור אחד \enLR{(DoS)} או ממקורות רבים \enLR{(DDoS).} & \textbf{\enLR{DoS / DDoS}} \\
\hline
הצפת שרת בבקשות \enLR{SYN} חצי-פתוחות עד שטבלת החיבורים מתמלאת. & \textbf{\enLR{SYN flood}} \\
\hline
שליחת קלט ארוך מהחוצץ, שגולש ודורס זיכרון ומריץ קוד תוקף. & \textbf{\enLR{Buffer overflow}} \\
\hline
ניסוי שיטתי של כל הסיסמאות האפשריות (או רשימת מילון) עד למציאה. & \textbf{\enLR{Brute force}} \\
\hline
שלב הסיור – איסוף מידע על המטרה \enLR{(whois}, סריקת פורטים, \enLR{sniffing).} & \textbf{\enLR{Reconnaissance}} \\
\hline
פגיעות שאין לה עדיין תיקון או שאינה ידועה ליצרן. & \textbf{\enLR{Zero-day}} \\
\hline
\end{xltabular}
\end{document}
